\section{Lean definitions, checked glue, and the completion path}
\label{sec:lean}

\subsection{What is compiled}

The archive contains a small reproducible Lake project pinned to Lean 4.32.0, Mathlib revision
\code{81a5d257c8e410db227a6665ed08f64fea08e997}, and the three unchanged ConciseIPRA modules \code{Syntax.lean}, \code{Deduction.lean}, and \code{PR.lean} from the source revision. The source hashes and upstream license are included.

\code{ObjectObservation.lean} compiles five generic inference combinators whose conclusions have the actual embedded \code{Prf PRAAx} type. They are not assertions about the standard natural numbers substituted for derivations. Their roles are:
\begin{enumerate}
\item \code{boundedInvariantAt}: use permitted quantifier-free induction and instantiate at a supplied numerical fuel term;
\item \code{loopProducesCertificate}: combine the terminal invariant with a supplied object-level extraction implication;
\item \code{observableOfIndependentCertificates}: apply a supplied stability implication to two derived certificates;
\item \code{observableFromLocalStability}: discharge two local certificate assumptions and apply the resulting implication;
\item \code{twoLoopsStableObservable}: combine two finite-loop derivations and a local observation comparison.
\end{enumerate}

Their axiom audit reports \code{propext} and \code{Quot.sound}, with no \code{sorryAx} or \code{Classical.choice}. These are Lean metatheoretic dependencies and do not add axioms to the embedded object calculus.

\code{SpectralBlueprint.lean} separately gives compiler-checked host definitions of finite rational matrices, sums, algorithms, certificates and target propositions. Definitions ending in \code{Target} name unproved propositions; they are neither axioms nor proved theorems. The structure \code{UnprovedMatrixObligations} lists them and has no supplied instance. Its rational implementation inherits \code{Classical.choice} transitively through Std, in addition to \code{propext} and \code{Quot.sound}. This host dependency is explicitly recorded. The file is not used as a derivation inside $\ipra$; the numerical reification below must have its own arithmetic proofs.

\subsection{The exact upstream interfaces}

The primitive recursive syntax is \code{PR : Nat -> Type}, generated by zero, successor, projection, composition and primitive recursion. Its term type is \code{PTm n}, generated by variables and application of a PR symbol. The step arguments for primitive recursion are ordered \emph{time, previous value, parameters}. Formula syntax and intuitionistic natural deduction are supplied by \code{Syntax.lean} and \code{Deduction.lean} \cite{prsource,deductionsource}.

For $\varphi$ over $n+1$ free variables, the existing rule \code{PRA.qfInd} takes
\[
 \varphi.\mathrm{QF},\qquad
 \Gamma\vdash_{\mathrm{PRAAx}}\varphi[0],\qquad
 \varphi::\mathrm{shift}(\Gamma)\vdash_{\mathrm{PRAAx}}\varphi[Sv_0]
\]
and returns $\Gamma\vdash_{\mathrm{PRAAx}}\forall v_0\,\varphi$.
The new \code{boundedInvariantAt} simply follows this with universal elimination at the explicit fuel term. The context may contain quantified input contracts; the induction formula itself must remain quantifier-free.

The final existing bridge is
\[
 (\Gamma\vdash_{\mathrm{PRAAx}}\varphi)
 \ \longleftrightarrow\
 \operatorname{NPrf}(\Gamma.\mathrm{map}\,\operatorname{NFm.ofPRA},
                            \operatorname{NFm.ofPRA}(\varphi)),
\]
named \code{pra_iff_native} \cite{nativesource}. Applying its forward direction is the last step after the desired embedded proof has been built. It is not a compiler from arbitrary Mathlib theorems to $\ipra$.

\subsection{The concrete host definitions}

The blueprint uses
\begin{lstlisting}
abbrev Vec (n : Nat) := Fin n -> Rat
abbrev Mat (n : Nat) := Fin n -> Fin n -> Rat

def sumQ : (n : Nat) -> (Fin n -> Rat) -> Rat
  | 0, _ => 0
  | n + 1, f => f 0 + sumQ n (fun i => f i.succ)
\end{lstlisting}
From this recursive sum it defines transpose, multiplication, diagonal matrices, reconstruction $QDQ^{\mathsf T}$, conjugation, squared Frobenius norm, off-diagonal energy, matrix-vector multiplication and quadratic forms. There is no real completion in these definitions.

Certificate data and validity are deliberately separated. \code{DiagonalData n} has a rational basis and vector of values. \code{DiagonalCertificate A eta w} states symmetry, nonnegative tolerance, both orthogonality identities, sortedness, exact agreement with the transformed diagonal, and the squared residual bound. \code{PositiveData} has a general rational factor and nonnegative weights, not an orthogonal factor; otherwise generic positive rational matrices would be excluded. \code{RootCertificate} asserts a positive root candidate and a squared residual. \code{ProjectorCertificate} asserts symmetry, idempotence and a commutator residual.

The file also defines the exact two-branch \code{pivotRotation}, plane rotations, finite first-maximum pivot selection, the frozen \code{jacobiStep}, primitive-recursive \code{jacobiRun}, the explicit \code{jacobiFuel}, and diagonal sorting. A harmless total fallback in the sorted-list lookup is eliminated by the usual sorting-length proof. The producer's mathematical correctness is the unproved proposition \code{DiagonalizationFuelTarget}, not inferred merely from type checking its definition.

\subsection{Numerical encodings for the object layer}

One explicit encoding is sufficient to make the path unambiguous. Use signed-integer codes $0\mapsto0$, positive $z\mapsto2z$, and negative $-z\mapsto2z-1$ for $z>0$. A convenient positive pairing is
\[
 \operatorname{pair}(a,b)=2^a(2b+1).
\]
Its projections use bounded search for the exponent of two and integer division. The code zero has fixed default projections. For a positive pair the second component is smaller than the pair code, and all projections are primitive recursive. Use pair-coded signed numerator and denominator-minus-one for rationals, with zero again given the default value $0/1$.

Encode a list by its length and iterated positive pairs, and a matrix by its dimension and a row-major list of length $n^2$. A validity checker verifies exactly the requested number of cells. Every raw operation is totalized on invalid codes. Cross-multiplication proves rational equality and order; normalization can be omitted from the logical encoding. For the computational bit bound use the reduced or common-denominator representation specified in Section~\ref{sec:grid}.

The fixed functions to reify as \code{PR} syntax include rational arithmetic; finite sums and products; array access/update; pivot scanning; plane-rotation update; state iteration; rational ceiling; final permutation; the positivity stack and its reconstruction; greatest-square bounded search; and threshold masks. The ordinary PR closure theorem applies to each fixed definition. A \code{PR} code for each function and a correctness relation to its host meaning are construction tasks, not supplied automatically by the Python program or Lean's termination checker.

\subsection{The proof obligations in dependency order}

\begin{longtable}{@{}P{0.25\textwidth}P{0.43\textwidth}P{0.26\textwidth}@{}}
\toprule
Target & Main proof and induction data & Required output \\
\midrule\endhead
Arithmetic representation & Signed order; positive denominators; cross-products; finite Boolean folds. Induct on numerical recursion counters with all parameters fixed. & Object equations and reflection lemmas for each fixed checker. \\
Finite matrix algebra & Expand bounded sums; prove reindexing, transpose and multiplication laws; finite Cauchy--Schwarz by the sum-of-squares identity. & Encoded identities and squared inequalities. \\
Pivot improvement and energy identity & Rational branch split; guarded denominator clearing; the two scalar inequalities and preservation of outside-pivot squared sums. & Local object proof for one finite update. \\
Finite diagonalization & The Booleanized invariant \eqref{eq:invariant}, exact orthogonality and conjugation, maximum-pivot scan, final sorting. & \code{Prf PRAAx} of producer correctness at the explicit fuel. \\
Rational intersection & Row-operation stack, pivot/free columns, back substitution, and nonzero free coordinate; at most $n$ stages. & A PR rational intersection witness with its finite checker proof. \\
Sorted spectral comparison & Use that witness, sorted coordinate bounds and finite Cauchy--Schwarz; clear the positive $x^{\mathsf T}x$ denominator. & Object cross-input comparison for arbitrary certificates. \\
Positivity targets & Bounded Schur elimination and reverse reconstruction, exact congruence invariant, negative-vector alternative. & Producer correctness for weighted positive factors or negative witnesses. \\
Root producer and mixed-basis comparison & Bounded square-grid search, exact orthogonal reconstruction, weighted Cauchy--Schwarz. & Root residual and comparison; then the finite uniqueness modulus. \\
Projector targets & Bits square to themselves; mixed-basis separation; positive gap matrix; finite row/column squared sums. & Projector identity, residual, gap margin, and cross-input comparison. \\
Name schemas & Instantiate the finite comparison at computed precisions, under fixed universal input-name contracts. & Arithmetic finite-observation theorems for every output precision. \\
Native transfer & Apply the already proved \code{pra_iff_native.mp}. & Native-core derivation of the translated formula. \\
\bottomrule
\end{longtable}

Every variable-size statement in the matrix rows is expressed by a checker on a finite code. For example, the orthogonality component is the Boolean product over $i,j<n$ of the equality tests for the corresponding finite dot products. The invariant does not quantify over all rational vectors or all future precisions. The universal quadratic-form theorem is established by proving the finite identity with its vector as a free numerical parameter and generalizing afterward. Similarly, a matrix name's universal regularity or positivity statement remains a fixed context assumption instantiated at needed numerical arguments.

For determinant multiplication, expand the two permutation sums. Terms indexed by a repeated intermediate column cancel in pairs under the transposition of the repeated positions; the remaining terms are indexed by permutations and factor into the two determinant sums. This is a bounded involution and reindexing argument. The principal-minor coefficient identity follows by choosing which diagonal $X$ terms occur in the permutation expansion. Its finite binomial and factorial counts have their own prefix-count invariants. These supply the elementary algebra required by Lemma~\ref{lem:coeff} without importing a semantic spectral theorem.

\subsection{What would count as a completed formalization}

A completed application must contain actual numerical PR definitions, object-level correctness proofs for the finite arithmetic kernels, the producer and comparison derivations, and the final instantiated theorem. Supplying an object formula with an axiom naming the desired theorem would not meet this criterion. Neither would proving only its semantic interpretation in Lean, even if that ambient proof is constructive.

The supplied compiled combinators already enforce the right interface: their missing local estimates are arguments of type \code{Prf PRAAx}, not arbitrary Lean propositions. They cannot silently absorb a host spectral theorem. The blueprint and reference programs make the remaining mathematical obligations concrete and testable, but the main matrix-specific embedded proofs remain to be constructed.
